\section{Market Clearing Conditions at Fixed Regime and Stationary Objects}





















































































































































































































































































































This phase derives the benchmark market-clearing block conditional on candidate incumbent measures and entrant masses. The purpose is to formalize qualified contract-manufacturing market clearing and labor market clearing under regime $M\in\{0,1\}$. This section does not add new benchmark mechanisms. It only states explicit regularity conditions that make the equilibrium equations well defined and solvable. The incumbent measures used here are treated as given inputs for the price-clearing block; Phase 7 later imposes stationary consistency on those measures.

\subsection{Benchmark Fidelity Check}

\paragraph{Step 1: benchmark objects used in this phase.}
The canonical market-clearing objects are:
\begin{itemize}
    \item $x_B^*(z;p_m,M)$ and $\lambda_B(z,p_m,M)$ from the type-$B$ static block;
    \item $m^*(z;p_m,w)$ from the type-$C$ static block;
    \item candidate incumbent measures $\mu_A,\mu_B,\mu_C$, treated as given finite Borel measures in this phase, and entrant masses $n_E^A,n_E^B,n_E^C$;
    \item type-specific labor-demand objects $\ell_A(z;w)$, $\ell_B(z;p_m,M)$, $\ell_C(z;p_m,w)$, and entrant setup labor coefficients $\ell_E^j$;
    \item exogenous labor supply $\bar L$.
\end{itemize}

\paragraph{Step 2: benchmark equations.}
Qualified contract-manufacturing demand and supply are
\begin{align}
D_m(p_m,M)
&=\int_Z x_B^*(z;p_m,M)\,\lambda_B(z,p_m,M)\,\mu_B(dz),
\label{eq:p6_Dm}\\
S_m(p_m,w)
&=\int_Z m^*(z;p_m,w)\,\mu_C(dz),
\label{eq:p6_Sm}
\end{align}
and market clearing requires
\begin{equation}
D_m(p_m,M)=S_m(p_m,w).
\label{eq:p6_pm_clear}
\end{equation}

The main manuscript later compresses the labor block into the equilibrium object $L_D(M)$. In the present fixed-price derivation, we therefore introduce the auxiliary fixed-price labor-demand map
\begin{equation}
\begin{aligned}
\mathfrak L_D(w,p_m,M)={}&\int_Z \ell_A(z;w)\,\mu_A(dz)
+\int_Z \ell_B(z;p_m,M)\,\mu_B(dz)\\
&+\int_Z \ell_C(z;p_m,w)\,\mu_C(dz)
+\sum_{j\in\{A,B,C\}}\ell_E^j n_E^j(M),
\end{aligned}
\label{eq:p6_LD_aux}
\end{equation}
and labor-market clearing requires
\begin{equation}
\mathfrak L_D(w,p_m,M)=\bar L.
\label{eq:p6_labor_clear}
\end{equation}

\paragraph{Step 3: benchmark-fidelity boundary.}
The contract-manufacturing block can be written directly from the benchmark objects above. The labor block can also be written faithfully, but the fixed-price notation $\mathfrak L_D(w,p_m,M)$ is auxiliary and is introduced only to avoid silently expanding the benchmark object $L_D(M)$. In addition, the current benchmark sources do not provide a fully primitive derivation of aggregate labor-demand continuity from lower-level labor policies, so the labor regularity conditions below are stated explicitly as supplemental closure assumptions.

\paragraph{Step 4: boundary discipline.}
This phase keeps the benchmark boundary: no government optimization, no DSGE closure, no bargaining/acquisition layer, and no extra state dimensions.

\subsection{Part 1: Qualified Contract-Manufacturing Market}

\paragraph{Assumptions for well-defined demand and supply.}
Fix regime $M$ and wage $w$.
\begin{assumption}[MC1: Integrability]
The mappings $z\mapsto x_B^*(z;p_m,M)\lambda_B(z,p_m,M)$ and $z\mapsto m^*(z;p_m,w)$ are Borel measurable and integrable under $\mu_B$ and $\mu_C$ for each admissible $p_m$.
\end{assumption}

Before stating the aggregate continuity assumption, it is useful to separate what is already available from Phase 2 and what is newly assumed here. Phase 2 derives continuity of $x_B^*(z;p_m,M)$ and $m^*(z;p_m,w)$ on compact price domains. Therefore, if $\lambda_B(z,p_m,M)$ is continuous in $p_m$ and the candidate measures $\mu_B,\mu_C$ are finite, dominated-convergence arguments imply continuity of the aggregate demand and supply maps below. MC2 records exactly that aggregate-level regularity conclusion because the market-clearing theorem uses the aggregate objects, not the pointwise policy functions.

\begin{assumption}[MC2: Continuity in $p_m$]
$D_m(p_m,M)$ and $S_m(p_m,w)$ are continuous in $p_m$ on a compact feasible interval $\mathcal P_m=[\underline p_m,\overline p_m]\subset\mathbb R_+$.
\end{assumption}

\begin{assumption}[MC3: Monotonicity]
$D_m(\cdot,M)$ is weakly decreasing and $S_m(\cdot,w)$ is weakly increasing on $\mathcal P_m$.
\end{assumption}

\begin{assumption}[MC4: Boundary crossing]
At the endpoints,
\[
D_m(\underline p_m,M)-S_m(\underline p_m,w)\ge 0,
\qquad
D_m(\overline p_m,M)-S_m(\overline p_m,w)\le 0.
\]
\end{assumption}

\begin{proposition}[Existence of qualified-manufacturing clearing price]
Under MC1--MC4, there exists at least one $p_m^*(w,M)\in\mathcal P_m$ satisfying
\[
D_m(p_m^*,M)=S_m(p_m^*,w).
\]
If, in addition, $D_m(\cdot,M)$ is strictly decreasing and $S_m(\cdot,w)$ is strictly increasing, the clearing price is unique.
\end{proposition}

\begin{proof}
Define excess demand
\[
\Delta_m(p_m;w,M):=D_m(p_m,M)-S_m(p_m,w).
\]
By MC2, $\Delta_m$ is continuous on compact interval $\mathcal P_m$. By MC4, $\Delta_m(\underline p_m;w,M)\ge 0$ and $\Delta_m(\overline p_m;w,M)\le 0$. The intermediate value theorem yields at least one $p_m^*$ with $\Delta_m(p_m^*;w,M)=0$.

For uniqueness, suppose $p_1<p_2$ both clear. Strict monotonicity implies
\[
\Delta_m(p_1;w,M)>\Delta_m(p_2;w,M),
\]
which contradicts $\Delta_m(p_1;w,M)=\Delta_m(p_2;w,M)=0$. Hence the solution is unique.
\end{proof}

\subsection{Part 2: Labor Market}

\paragraph{Assumptions for labor-demand closure.}
Fix regime $M$ and a candidate $p_m$.
Unlike the contract-manufacturing block, the current benchmark sources do not provide a fully primitive derivation of aggregate labor-demand continuity from lower-level labor policy functions. For that reason, LC1--LC3 are best read as supplemental closure assumptions on the auxiliary fixed-price labor-demand map rather than as already-proved consequences of earlier phases.
\begin{assumption}[LC1: Integrability and continuity]
$\mathfrak L_D(w,p_m,M)$ in \eqref{eq:p6_LD_aux} is finite and continuous in $w$ on compact wage interval $\mathcal W=[\underline w,\overline w]\subset\mathbb R_+$.
\end{assumption}

\begin{assumption}[LC2: Monotonicity in wage]
$\mathfrak L_D(\cdot,p_m,M)$ is weakly decreasing on $\mathcal W$.
\end{assumption}

\begin{assumption}[LC3: Boundary crossing]
\[
\mathfrak L_D(\underline w,p_m,M)\ge \bar L,
\qquad
\mathfrak L_D(\overline w,p_m,M)\le \bar L.
\]
\end{assumption}

\begin{proposition}[Existence of labor-clearing wage]
Under LC1--LC3, for each fixed $(p_m,M)$ there exists at least one wage $w^*(p_m,M)\in\mathcal W$ such that
\[
\mathfrak L_D(w^*,p_m,M)=\bar L.
\]
If $\mathfrak L_D(\cdot,p_m,M)$ is strictly decreasing, the wage is unique.
\end{proposition}

\begin{proof}
Define excess labor demand
\[
\Delta_L(w;p_m,M):=\mathfrak L_D(w,p_m,M)-\bar L.
\]
By LC1, $\Delta_L$ is continuous on $\mathcal W$. By LC3, $\Delta_L(\underline w;p_m,M)\ge 0$ and $\Delta_L(\overline w;p_m,M)\le 0$. Thus IVT implies existence of at least one root $w^*$.

If $\mathfrak L_D(\cdot,p_m,M)$ is strictly decreasing, then $\Delta_L(\cdot;p_m,M)$ is strictly decreasing and can cross zero at most once, implying uniqueness.
\end{proof}

\subsection{Part 3: Joint Price-Clearing Mapping (Interface Statement)}

The two markets are coupled through $(p_m,w)$. This phase establishes one-dimensional existence results conditional on the other price. The full stationary-equilibrium existence argument can now treat the pair
\[
\begin{cases}
D_m(p_m,M)-S_m(p_m,w)=0,\\
\mathfrak L_D(w,p_m,M)-\bar L=0,
\end{cases}
\]
as a joint fixed-point system on $\mathcal P_m\times\mathcal W$.

\paragraph{Why this is enough for Phase 6.}
Phase 6 does not need to re-solve Bellman objects or entry logic. It only has to provide a clean and verifiable market-clearing closure interface for later stationary-measure and full-equilibrium phases.

\subsection{What did we just do, and why does it matter?}

This phase separated the two price-clearing conditions that sit between the firm's dynamic problem and the later equilibrium block. Throughout the phase, the incumbent measures are treated as candidate inputs rather than as already-derived stationary objects; Phase 7 later supplies the stationary-consistency argument. On the contract-manufacturing side, the object of interest is the scalar equation
\[
D_m(p_m,M)=S_m(p_m,w),
\]
which says that the benchmark demand for qualified contract-manufacturing services must equal the benchmark supply generated by type-$C$ incumbents at the candidate prices. On the labor side, the object of interest is
\[
\mathfrak L_D(w,p_m,M)=\bar L,
\]
which is the fixed-price auxiliary representation of the benchmark labor-clearing object $L_D(M)=\bar L$. By writing each market-clearing condition in one dimension and proving existence separately, the phase makes clear which assumptions do the work: continuity gives a root-finding problem that is well posed, monotonicity delivers economic discipline, and endpoint inequalities let the Intermediate Value Theorem apply.

This matters because later equilibrium arguments cannot simply assert that prices exist. They need explicit scalar conditions that can be checked, combined, and eventually embedded in a finite-dimensional outer map. Phase 6 therefore plays an interface role. It does not redesign the model, re-open the firm's optimization problem, or solve the full equilibrium. Instead, it takes candidate measures and entrant masses as given and shows exactly how the benchmark market-clearing block can be written in a theorem-ready form that later phases can use.

\subsection*{End-of-Section Recap}

\begin{itemize}
    \item What has been established mathematically: the qualified contract-manufacturing market-clearing equation and the labor-market clearing equation are both well defined, and each admits an existence result under explicit continuity, monotonicity, and boundary-crossing conditions.
    \item What assumptions were used: candidate incumbent measures and entrant masses treated as given in this phase, together with MC1--MC4 for the contract-manufacturing block and LC1--LC3 for the auxiliary fixed-price labor block.
    \item What remains to be derived next: the next phase must impose stationary consistency on the candidate incumbent measures and then combine those stationary objects with the market-clearing conditions.
\end{itemize}
